%!TEX root = ../main.tex
% Abstract in English (~250 words). Keywords are set in main.tex
% with \keywords{...}. Must be a faithful translation of the Resumen.

\begin{abstract}

Injection vulnerabilities, such as SQL or command injection, arise when
user-controlled data reaches a sensitive operation without protection.
Tools that look for them through pattern matching tend to flag every
sensitive operation even when the data is protected, and the resulting
volume of false alarms leads development teams to disregard the results.
GraphSAST is a static source code security analyzer that represents the
program as a data flow graph and, through taint analysis, follows the
path of untrusted data (sources) to sensitive operations (sinks),
discarding the paths that go through sanitization or validation. It
analyzes applications written in TypeScript and JavaScript, tracking data
across functions and files, as well as Python, Java, PHP and Go. It is
used from the command line, which can be integrated into continuous
integration, and from a web application that analyzes the code in the
browser, without sending it to a server, and shows the path of each
finding as a graph, together with its explanation and its fix.

On seven public repositories with 34 documented vulnerabilities,
precision was 0.41 against 0.10 for a regular expression baseline, with
a recall of 0.82 against 0.88. On the OWASP Benchmark for Java it
detected all 777 vulnerable cases with a false positive rate of 0.26,
and in a field test with 50 GitHub projects the web version and the
command line produced the same results. The hypothesis that the graph
makes the risk easier to understand than a text report remains
preliminary: the user validation gathered three participants.

\end{abstract}
